\FloatBarrier
\section{Communication between Qubits}
\label{app:qubit-communication}

\label{sec:interqubit-communication}

This appendix expands the final question in
Section~\ref{sec:platform-array-selection}: how qubits without a direct
interaction link can participate in a gate.

A circuit can request a gate between qubits whose vertices share
no interaction edge. The missing edge is a physical problem, not
merely a change of gate notation. We can bring the two quantum
states next to one another with SWAP gates, transport an electron
carrying its spin, or couple stationary spins through another
quantum system. These choices move different objects and have
different conditions for preserving an unknown state.

\subsection{Moving a state with SWAP gates}
\label{sec:processor-swap-routing}

An exchange pulse with the calibrated SWAP area exchanges the
spin states at neighbouring occupied dots
[Eq.~\eqref{eq:partial-swap}]. For example,
\begin{equation}
 \ket{\psi}_1\ket0_2\ket0_3
 \xrightarrow{\mathrm{SWAP}_{12}}
 \ket0_1\ket{\psi}_2\ket0_3
 \xrightarrow{\mathrm{SWAP}_{23}}
 \ket0_1\ket0_2\ket{\psi}_3.
 \label{eq:processor-swap-transfer}
\end{equation}
Each dot still holds an electron. The state \(\ket{\psi}\),
which may be unknown or entangled with other qubits, changes
its site because SWAP permutes spin states. The displayed
\(\ket0\) states make the path easy to read; intermediate
qubits need not be initialized for the permutation to work.
Their own states are permuted as well, and the control program
must track the resulting placement.

This permutation builds a gate between distant states from a
gate available only on an edge. In the three-site chain let
\(S_{23}\) be SWAP on sites 2 and 3, and let \(U_{12}\) be the
desired two-qubit operation on edge \((1,2)\). Then
\begin{equation}
 U_{13}=S_{23}U_{12}S_{23}.
 \label{eq:processor-routed-gate}
\end{equation}
The rightmost SWAP first moves the state originally at site 3
to site 2; \(U_{12}\) acts on that state and the one at site 1;
the leftmost SWAP restores both locations. The identity holds
on arbitrary three-qubit inputs, so routing neither measures
nor copies the transported state.

For a chosen path of \(d\) edges, bringing one state beside
the other takes \(d-1\) SWAPs; restoring the placement takes
another \(d-1\). If all are serial and have duration
\(t_{\mathrm{SWAP}}\), a routed gate of duration \(t_U\) takes
\[
 T_{\mathrm{route}}=2(d-1)t_{\mathrm{SWAP}}+t_U.
\]
In an illustrative model where each SWAP independently fails
with probability \(p_S\) and the central gate fails with
probability \(p_U\), the probability of at least one fault is
\[
 p_{\mathrm{fault}}
 =1-(1-p_S)^{2(d-1)}(1-p_U)
 \simeq2(d-1)p_S+p_U.
\]
This stochastic estimate is not a rule for adding coherent
gate errors or average infidelities. If later operations can
use the new placement, the restoring SWAPs may be omitted,
provided the program tracks every moved state. For a
singlet--triplet or exchange-only encoding, swapping one
physical spin generally does not swap the whole logical
qubit; the logical-block operation must be constructed and
checked separately.

\subsection{Moving an electron while preserving its spin}
\label{sec:processor-shuttling}

SWAP moves a spin state through occupied sites. Shuttling instead
moves the electron that carries the spin into another dot. For one
electron in dots \(A\) and \(B\), its position is a charge
two-level system with
\(\tau_z=\ket A\bra A-\ket B\bra B\) and
\(\tau_x=\ket A\bra B+\ket B\bra A\). A minimal model is
\begin{equation}
 H(t)=
 \left[\frac{\varepsilon(t)}2\tau_z+t_c\tau_x\right]\otimes I
 -\frac12\left[E_{Z,A}\ket A\bra A+
              E_{Z,B}\ket B\bra B\right]\otimes Z,
 \label{eq:processor-shuttle-hamiltonian}
\end{equation}
where \(\varepsilon=\epsilon_A-\epsilon_B\), \(t_c>0\) is
spin-conserving tunneling, and \(Z\ket0=\ket0\) for
\(\ket0=\ket{\downarrow}\). A sweep from negative to positive
\(\varepsilon\) changes the charge ground state from mostly
\(\ket A\) to mostly \(\ket B\). Successful transport of the
electron does not, by itself, establish that its spin state
survived the trip.

First neglect the differential Zeeman energy when finding that
charge ground state. Its gap to the excited charge state and its
probability of occupying \(A\) are
\[
 \Delta_c=\sqrt{\varepsilon^2+4t_c^2},\qquad
 n_A=\frac12\left(1-\frac{\varepsilon}{\Delta_c}\right).
\]
Thus \(n_A\) changes continuously from one to zero and equals
one half at the avoided crossing. To follow the ground branch,
the local adiabatic condition is
\begin{equation}
 \frac{\hbar t_c|\dot\varepsilon|}
      {(\varepsilon^2+4t_c^2)^{3/2}}\ll1.
 \label{eq:processor-shuttle-speed}
\end{equation}
It has the form derived for a charge ramp in
Eq.~\eqref{eq:calibration-charge-ramp-condition}, with the
single-electron coupling \(t_c\). The crossing is the
restrictive part of a linear sweep because the gap is smallest
there.

One way to allocate more time near the crossing is to hold the
dimensionless quantity in Eq.~\eqref{eq:processor-shuttle-speed}
at a chosen small value \(\eta_{\mathrm{ad}}\) over the main
part of the ramp. Then
\(\dot\varepsilon=
\eta_{\mathrm{ad}}(\varepsilon^2+4t_c^2)^{3/2}/(\hbar t_c)\).
Integrating between equal and opposite endpoint detunings gives
\begin{equation}
 T=\frac{\hbar\varepsilon_{\max}}
 {2\eta_{\mathrm{ad}}t_c
                 \sqrt{\varepsilon_{\max}^2+4t_c^2}}.
 \label{eq:processor-local-adiabatic-time}
\end{equation}
Smooth connections to the endpoints alter the total pulse time.
Moreover, \(\eta_{\mathrm{ad}}\) is a local condition, not a
prediction of the full transfer error. Charge readout after the
sweep tests whether the electron reached \(B\).

Even perfect charge following can rotate the spin. Provided
\(|E_{Z,A}-E_{Z,B}|\ll\Delta_c\) throughout the sweep, projection
onto the charge ground state gives
\[
 H_{\mathrm{spin}}(t)=-\frac12E_Z(t)Z,\qquad
 E_Z(t)=n_A(t)E_{Z,A}+[1-n_A(t)]E_{Z,B}.
\]
The electron therefore samples both local Zeeman energies as
it crosses. If spin flips and charge excitation are negligible,
\begin{equation}
 \ket A\otimes(a\ket0+b\ket1)
 \longrightarrow
 \ket B\otimes R_z(-\varphi)(a\ket0+b\ket1),
 \qquad
 \varphi=\frac1\hbar\int_0^T E_Z(t)\,dt,
 \label{eq:processor-shuttle-phase}
\end{equation}
up to a common orbital phase. This phase is written in the
laboratory frame. Relative to a rotating control reference of
angular frequency \(\omega\), the measured phase is instead
\(\int_0^T[E_Z(t)/\hbar-\omega]\,dt\). A reproducible phase can
be included in the control frame; an uncertain phase cannot.

This distinction requires more than an arrival measurement.
For a transferred \(\ket+\), \(X\) and \(Y\) analyses in the
chosen frame yield \(\langle X\rangle=\cos\varphi\) and
\(\langle Y\rangle=-\sin\varphi\). If the phase varies between
shots, the combined contrast becomes
\[
 \sqrt{\langle X\rangle^2+\langle Y\rangle^2}
       =|\mathbb E[e^{-i\varphi}]|.
\]
A fixed phase rotates the vector without reducing its length;
phase fluctuations shorten it. Basis-state transfers also test
spin flips. Charge arrival, spin preservation conditioned on
arrival, and the reproducibility of the accumulated phase are
therefore separate properties of a shuttle. Silicon experiments
have tested coherent spin transport with these distinct
observables~\cite{Yoneda2021Transport}.

\subsection{Coupling through an electromagnetic mode}
\label{sec:processor-mode-coupling}

A third route leaves both electrons at their dots. Each spin
couples to the same resonator mode, and an excitation can pass
between the spins through that mode. Its field is a quantum
degree of freedom: unlike a prescribed microwave drive, it can
receive an excitation and later return it.

One mechanism converts the resonator's electric field into
spin coupling through a magnetic-field gradient. Let \(a\)
annihilate a photon of angular frequency \(\omega_c\), and let
\(E_{\mathrm{zpf}}(a+a^\dagger)\) be the electric field along the
electron's displacement. In a harmonic dot with orbital
frequency \(\omega_{\mathrm{orb}}\gg\omega_c,\omega_i\),
\[
 x\simeq-\frac{eE_{\mathrm{zpf}}}
                  {m^*\omega_{\mathrm{orb}}^2}(a+a^\dagger).
\]
A transverse field gradient \(b_x=\partial B_x/\partial x\)
then turns the displacement into the spin term
\(g_e\mu_B b_x xX/2\), giving an angular-frequency coupling
of magnitude
\[
 g_s=\frac{|g_e\mu_B b_x eE_{\mathrm{zpf}}|}
                {2\hbar m^*\omega_{\mathrm{orb}}^2}.
\]
Here \(g_e\) is the electron Land\'e factor. This expression
explains the route from an electric field to a magnetic spin
drive. Actual charge eigenstates can change the coupling, so
\(g_s\) is determined for the device rather than inferred from
this harmonic estimate alone. Spin--photon coupling has been
observed in silicon~\cite{Samkharadze2018}.

With measured couplings \(g_1,g_2\), the
excitation-conserving model is
\begin{equation}
 \frac H\hbar=\omega_c a^\dagger a+
   \sum_{i=1}^2\omega_i\ket1_i\bra1+
   \sum_{i=1}^2g_i(a\sigma_i^++a^\dagger\sigma_i^-),
 \qquad \sigma_i^+=\ket1_i\bra0.
 \label{eq:processor-spin-mode-model}
\end{equation}
The rotating-wave approximation requires
\(g_i,|\omega_i-\omega_c|\ll\omega_i+\omega_c\).
When one spin is resonant with an initially empty mode, an
initial \(\ket{1,0_{\mathrm{ph}}}\) oscillates with
\(P_1(t)=\cos^2(g_i t)\) if losses are negligible during the
hold. The mode is then populated during the exchange.

A detuned mode can instead act mainly as a virtual mediator.
For equal bare spin frequencies \(\omega_q\), put
\(\Delta=\omega_q-\omega_c\) and assume
\(|\Delta|\gg g_1,g_2\). In the one-excitation basis
\(\{\ket{10,0_{\mathrm{ph}}},\ket{01,0_{\mathrm{ph}}},
\ket{00,1_{\mathrm{ph}}}\}\), after subtracting
\(\hbar\omega_q\),
\[
 \frac{H_1}{\hbar}=
 \begin{pmatrix}0&0&g_1\\0&0&g_2\\g_1&g_2&-\Delta\end{pmatrix}.
\]
The mode amplitude follows the two spin amplitudes
approximately as \(c_c=(g_1c_1+g_2c_2)/\Delta\) on their slow
time scale. Substitution gives
\begin{equation}
 \frac{H_{\mathrm{eff},1}}{\hbar}
 =\frac1\Delta
  \begin{pmatrix}g_1^2&g_1g_2\\g_1g_2&g_2^2\end{pmatrix}.
 \label{eq:processor-mediated-exchange}
\end{equation}
Its diagonal shifts must be included when matching the two
spin frequencies; the off-diagonal interaction has magnitude
\(|g_{\mathrm{med}}|=|g_1g_2/\Delta|\). With the shifts
compensated, an initial \(\ket{10}\) evolves to
\(\cos(g_{\mathrm{med}}t)\ket{10}
-i\sin(g_{\mathrm{med}}t)\ket{01}\), apart from a common
phase. The equal-weight entangling time is
\(\pi/(4|g_{\mathrm{med}}|)\), and full excitation transfer
takes \(\pi/(2|g_{\mathrm{med}}|)\). Increasing \(|\Delta|\)
reduces the mode occupation, of order \((g_i/\Delta)^2\),
while lengthening both times. Whether the remote operation is
useful depends on this tradeoff and on spin decoherence and
resonator loss over the full gate.
